\section{The equality case}\label{sec:equality}

\Cref{prop:terminal-equality} identifies a period of cost $\Aoct$ with the octagon's trajectory,
but only after three choices that the proof of the lower bound made along the way: the
normalization of the disk to Reinhardt's normal position, the rebasing of its trajectory at a
switch, and a relabeling of the vertices of the control simplex by a power of $R$. To classify the
disks of density $\doct$ these choices must be undone, and the trajectory to which the terminal
equality is applied must be the realization of the given disk; another minimizer of the same
cost would say nothing about it. This section carries the identification back to the disk
(\cref{sec:equality-rigidity}), proves the converse (\cref{sec:equality-converse}), identifies
Hales's octagon $\KH$ with Reinhardt's $\Koct$ (\cref{sec:equality-construction}). These results
complete the proofs of \cref{thm:main}(a) and~(c) in \cref{sec:equality-proofs}.

\subsection{Rigidity}\label{sec:equality-rigidity}

Let $s_0,\dots,s_5$ be the sixth roots of unity, so that $Rs_j=s_{j+1}$. The realization of a worst
disk (\cref{thm:equivalence}) keeps the link between its frame and its boundary.

\begin{lemma}[Realization with its contact link]\label{lem:eq-realization}
Let $K$ be a worst disk. There are an invertible linear map $e$ of the plane, a balanced pair
$\sigma$ of $eK$ in circle representation and a global minimizer $(T,u,g,X)$ of the control
problem such that $eK$ is a worst disk, $\delta(K)=\cost(T,X)/\sqrt{12}$, and for every $j$ and
every $t\in\R$
\[
 \sigma_j(t)=g(t)\,s_j,\qquad \sigma_j(t+T)=\sigma_{j+1}(t).
\]
\end{lemma}

The first identity is the definition of the frame of the realization, which is defined on all
of $\R$; the second is the cyclic shift of the contacts by one sixth of the period. Together they
force the frame to be twisted on all of $\R$, beyond the period $[0,T]$ on which the control
problem is posed.

\begin{lemma}[The frame twist]\label{lem:eq-frame-twist}
If $\sigma_j(t)=g(t)s_j$ and $\sigma_j(t+T)=\sigma_{j+1}(t)$ for all $j$ and $t$, then
$g(t+T)=g(t)R$ for every $t\in\R$.
\end{lemma}

\begin{proof}
$g(t+T)s_j=\sigma_j(t+T)=\sigma_{j+1}(t)=g(t)s_{j+1}=g(t)Rs_j$, and $s_0,s_1$ span the plane.
\end{proof}

The last ingredient recovers a disk from its boundary.

\begin{lemma}\label{lem:eq-frontier}
Two compact convex sets with nonempty interior, both symmetric about the origin, that have the
same frontier are equal.
\end{lemma}

\begin{proof}
Let $A,B$ be the two sets. The interior of $B$ is connected, contains the origin, and does not
meet the common frontier; it is covered by the disjoint open sets $\inter A$ and
$\R^2\setminus\overline A$, so it lies in $\inter A$, which contains the origin. By symmetry the
interiors are equal, and a convex body is the closure of its interior.
\end{proof}

\begin{proof}[Proof of \cref{thm:rigidity}]
Let $K$ be a worst disk, and take $e$, $\sigma$ and $(T,u,g,X)$ as in \cref{lem:eq-realization}.
By \cref{lem:eq-frame-twist}, $g(t+T)=g(t)R$ on $\R$.

\emph{The cost is $\Aoct$.} Since $K$ is a worst disk and $\delta(\KH)=\doct$
(\cref{prop:octagon-trajectory}), $\delta(K)\le\doct$, and part~(b) of \cref{thm:main} gives
$\delta(K)\ge\doct$. Hence $\cost(T,X)=\sqrt{12}\,\delta(K)=\Aoct$.

\emph{The terminal equality.} By \cref{thm:finite-normal}, the minimizer rebased at some
$s\in[0,T)$ is a normal vertex period $P$ with the same terminal time and cost, whose state stays
in $\hstarstar$; its frame is $t\mapsto g(s)^{-1}g(s+t)$, because $g$ is already twisted. By
\cref{thm:at-most-five}, $P$ has at most five pieces, and by \cref{prop:terminal-equality},
$T=\tf$ and there is $r$ with $R^{-r}g(s)^{-1}g(s+t)R^{r}=g_{\mathrm{oct}}(t)$ for $t\in[0,T]$.
With $M=g(s)R^r$ this reads
\[
 g(s+t)=M\,g_{\mathrm{oct}}(t)\,R^{-r}\qquad(0\le t<T).
\]
Both sides satisfy $F(t+T)=F(t)R$: the left side by the twist of $g$, the right side because
$g_{\mathrm{oct}}(t+\tf)=g_{\mathrm{oct}}(t)R$ and $R$ commutes with $R^{-r}$. Two curves with the
same twist that agree on $[0,T)$ agree everywhere, so the identity holds for all $t\in\R$.

\emph{The arcs.} Applying it to $s_0$ and using $R^{-r}s_0=s_{-r}$,
\[
 \sigma_0(s+t)=M\,g_{\mathrm{oct}}(t)\,s_{-r}=M\,\sigma^{\mathrm{oct}}_{-r}(t)\qquad(t\in\R),
\]
where $\sigma^{\mathrm{oct}}_j(t)=g_{\mathrm{oct}}(t)s_j$ are the contact curves of the octagon.
In a balanced pair each contact curve traces the whole frontier once per period, so the frontier
of $eK$ is the range of $\sigma_0$, and the frontier of $\KH$ is the range of
$\sigma^{\mathrm{oct}}_{-r}$. Hence $\partial(eK)=M(\partial\KH)=\partial(M\KH)$. Both $eK$ and
$M\KH$ are convex disks symmetric about the origin, so $eK=M\KH$ by \cref{lem:eq-frontier}, and
$K=e^{-1}M(\KH)$.
\end{proof}

The proof never compares $K$ with a minimizer other than its own realization. It runs the
reduction of \crefrange{sec:reduction}{sec:terminal} a second time, on that realization, and uses
the lower bound only to know that its cost is $\Aoct$.

\subsection{The converse}\label{sec:equality-converse}

The converse is the affine invariance of the greatest packing density.

\begin{proof}[Proof of \cref{prop:converse}]
A translate, and more generally a congruent copy, of $K$ has the same packings by congruent copies,
so its greatest packing density is the same. By
Fejes T\'oth's theorem (\cref{thm:fejes-toth}), $\delta(K)$ is the greatest density
$\operatorname{area}(K)/\det\Lambda$ of a lattice packing $K+\Lambda$. If $\det A>0$, the
lattice packings of $K$ and of $AK$ correspond through $\Lambda\mapsto A\Lambda$, and area and
covolume are both multiplied by $\det A$, so the lattice densities agree. If $\det A<0$, write
$A=(AF)F$ with the reflection $F(x_0,x_1)=(x_0,-x_1)$: then $FK$ is a congruent copy of $K$, and
$AF$ preserves orientation. The last statement follows from $\delta(\KH)=\doct$
(\cref{prop:octagon-trajectory}).
\end{proof}

\subsection{Reinhardt's construction}\label{sec:equality-construction}

\Cref{thm:main} is stated for Reinhardt's octagon $\Koct$ (\cref{def:smoothed-octagon}), while
the proof works with $\KH$, the region swept from the origin by the contact curve
$\sigma^{\mathrm{oct}}_0$ of the octagon's trajectory over one period $[0,6\tf]$. By
\cref{prop:octagon-trajectory}, $\KH$ is also the convex hull of the six arcs
$\sigma^{\mathrm{oct}}_j([0,\tf])$. The two octagons differ by a vertical stretch of about $5.5\%$
(\cref{fig:octagon}).

\begin{proof}[Proof of \cref{prop:reinhardt-construction}]
Write $y_0^2=(2\sqrt2-1)/3$ and $W=\exp(\sqrt3\,y_0 s)$ for $s\in[0,\tsw]$, so that $W$ runs over
$[1,\sqrt2]$; recall that $\sigma^{\mathrm{oct}}_0(0)=s_0=(1,0)$.

\emph{One eighth of the boundary.} On the first switching interval the octagon's frame
$g_0(s)$ is the flow of the control $\ee_2$ from $\Phi(0,y_0)$; conjugated by $A$ it acts by
\[
 Ag_0(s)A^{-1}=\begin{pmatrix}1&-\frac{W-1}{2\sqrt2-1}\\[0.5ex]
 1-W^{-1}&1-\frac{W-2+W^{-1}}{2\sqrt2-1}\end{pmatrix},
\]
and $A\,g_0(\tsw)R^{-1}A^{-1}$ is the rotation by $-\pi/4$; both are direct computations from
the closed form of the flow of $\ee_2$. Since $g_{\mathrm{oct}}(t+\tsw)=g_0(\tsw)R^{-1}\,g_{\mathrm{oct}}(t)\,R$,
the curve $A\sigma^{\mathrm{oct}}_0$ is, on the first three switching intervals,
\begin{align*}
 &(1,\ 1-W^{-1}) &&\text{along side $0$},\\
 &\Bigl(\sqrt2+1-\tfrac{\sqrt2}{2}(W+W^{-1}),\ 1-\tfrac{\sqrt2}{2W}\Bigr)
  &&\text{around corner $0$},\\
 &\Bigl(\sqrt2-\tfrac W2,\ \tfrac W2\Bigr) &&\text{along side $1$}.
\end{align*}
The first piece runs from the midpoint $(1,0)$ of side~$0$ to the point $(1,1-\sqrt2/2)$, and the
third from $(\sqrt2-\tfrac12,\tfrac12)$ to the midpoint $(\sqrt2/2,\sqrt2/2)$ of side~$1$.
With $d_k$ the distance to the $k$th side line, the middle piece has $d_7=W/2$ and
$d_2=\sqrt2/(2W)$, so it lies on the hyperbola $d_7d_2=\sqrt2/4$; it starts at
$(1,1-\sqrt2/2)$ and ends at $(\sqrt2-\tfrac12,\tfrac12)$, the points where the first and the
third pieces end and start. The control on these three intervals is $\ee_2$, $\ee_0$, $\ee_1$,
and only the second bends the curve.

\emph{The rest.} Three switching intervals later the frame relation gives
$\sigma^{\mathrm{oct}}_0(t+3\tsw)=-h^3\sigma^{\mathrm{oct}}_0(t)$ with $h=g_0(\tsw)R^{-1}$, and
$-Ah^3A^{-1}$ is the rotation by $\pi/4$. Over the period
$6\tf=24\tsw$ the curve $A\sigma^{\mathrm{oct}}_0$ therefore traces eight rotated copies of the
three pieces, once around.

\emph{The two inclusions.} Every point $q$ of the three pieces lies in $\Koct$, and $\mu q$ does
not for $\mu>1$: on the sides this is the inequality $d_k\ge0$, and on the corner it is the
hyperbola condition, which strictly fails past the arc. Since $\Koct$ is invariant under the
rotation by $\pi/4$, the same holds for every point $A\sigma^{\mathrm{oct}}_0(t)$. The set $\Koct$ is star-shaped about $0$, so $A\KH\subseteq\Koct$.
Conversely, a point $p\ne0$ of $\Koct$ is a positive multiple $\lambda q$ of some
$q=A\sigma^{\mathrm{oct}}_0(t)$, by the intermediate value theorem and
$\sigma^{\mathrm{oct}}_0(t+3\tf)=-\sigma^{\mathrm{oct}}_0(t)$, and $\lambda\le1$ because
$\mu q\notin\Koct$ for $\mu>1$. Hence $p\in A\KH$.
\end{proof}

Each side of $\Koct$ keeps a straight middle part of length $2-\sqrt2$, between the points of
tangency $(1,\pm(1-\sqrt2/2))$ on side~$0$. The proposition transfers the facts known for $\KH$ to
$\Koct$; in particular, the affine images of the two octagons are the same sets.
